\section{Background}\label{sec:background}

%{\bf Preallocation WAL files in database systems.}
%Database systems frequently benefit from preallocation files to a fixed size prior to use, establishing a stable file structure that mitigates runtime penalties associated with dynamic block allocations and file system journaling. For example, OceanBase, a rising-star database system, initializes its data and WAL files by preallocating tens of gigabytes by default. Similarly, in RocksDB, log and data (SSTable) files are uniformly sized at 64MB by default~\cite{DB:RocksDB:TOS-2021}, facilitating straightforward preallocation as an optimization strategy. Furthermore, MySQL's InnoDB engine similarly adopts preallocation for its redo log architecture\cite{mysql-wl4925}. The engine preallocates sequential log files at server initialization. This preallocation strategy minimizes the overhead of block management and file system journaling, thereby enhancing overall system efficiency.

{\bf  Databases logging.} %Logging is a technique that many
Many databases rely on logging for consistency and durability.
They  persistently write  new or updated data, or operations, in a log file  for recovery in case of power outage or kernel panic
\cite{DB:OceanBase:VLDB-2022,lee2015waldio,mysql2013mysql}.
Log files are hardly read except for recovery, so they can be operated in the direct I/O mode, with the OS's page cache bypassed.
 Moreover, some databases, like the rising-star OceanBase~\cite{DB:OceanBase:VLDB-2022},
 %rising-star databases such as OceanBase ,
   initialize
  log files by preallocating tens of gigabytes by default. 
  % Mature database engines such as SQLite can also be tuned to preallocate log files\cite{lee2015waldio}.




\textbf{File I/O Path.}
%When an application issues a file operation,
%A file operation that an application (e.g., database) issues %, such as a file write or fdatasync,
%When an application issues a file operation, the I/O request 
%it 
A file operation %issued by an application
traverses a complex  stack   of multiple software layers 
before reaching the storage device. % (e.g., NVMe SSD).
The top    {\em syscall interface} %at which syscalls such as file write and fsync 
  transitions syscalls such as file write and fsync to the OS kernel  with parameter validation.
The {\em virtual file system} ({VFS}) then performs inode lookup, interacts with the OS page cache, 
manages file position tracking, and so on.%~\cite{love2010linux}.
%The file operation is next passed to 
VFS passes the  operation to
a {\em specific file system}, e.g., %such as the prevalent 
Ext4. % studied in this paper.
Ext4
maps   in-file offsets to device block numbers and   composes block I/O (bio) requests.
For a write targeting unallocated space, Ext4 allocates   blocks on demand
and updates the per-file mappings. These metadata changes will trigger
{\em journaling} for file system consistency, and Ext4 commits modified metadata
%are committed 
to a  {\em journal}
before in-place update~\cite{fsync:FastCommit:ATC-2024,fsync:iJournaling:ATC-2017}.
The {\em bio layer} and {\em I/O scheduler}  manage, merge, and schedule bio requests where necessary. 
%bio requests, merge requests , and schedule 
%requests
%~\cite{love2010linux}.
The {\em device driver} translates bio requests to device-specific commands. It finally collects completion or fail signals from the device and
notifies upper~layers.



\textbf{Software Tax.}	
Software layers impose heavy tax on the critical path of serving  logging writes.
%When traversing these  layers, 	software overhead will occur on the critical path.
Particularly, locating a target block for each write involves a lookup
in the file's extent tree, by which Ext4 stores the offset-to-block mappings per file. 
%The block allocation triggers %, which
The journaling caused by allocating blocks and updating  per-file mappings 
entails severe write amplification with multiple writes performed to  the journal and log file \cite{fsync:iJournaling:ATC-2017,fsync:FastCommit:ATC-2024}.



\begin{comment}
 

{\bf Software tax for storage I/Os.}
%Files are the primary format for organizing data.
When an application-level program like OceanBase calls for a file operation,
I/O request traverses multiple layers of the de facto storage stack, including the Virtual File System (VFS), file system (e.g., Ext4), page cache, bio layer, and device driver, until it reaches the device.
They mainly carry on two tasks. % prior to actual I/Os.
One is to locate the target disk blocks. This is done by the kernel file system that looks up in the per-file mapping with an in-file offset. Ext4 use an extent tree to record the mapping for each file.
Searching in one such on-disk tree is non-trivial. The other task is to check access permissions.
Linux uses two protection models for file operations: the common Discretionary Access Control (DAC) and the less common Mandatory Access Control (MAC). DAC checks access
permission only at opening a file, whereas MAC happens at each I/O request.

\end{comment}



% huyp: have been mentioned in the intro
\begin{comment}
   Emerging memory technologies, like Intel 3D XPoint \cite{intel3dxpoint} and low-latency NAND flash \cite{samsungznand2024}, have paved the way for NVMe SSDs that offer $\mu$s-level latency and millions of I/O operations per second (IOPS).
With increasingly faster devices, the tax imposed by software layers turns to rise in the overall time cost per I/O request. 

{\bf SOTA solutions.}
Researchers have studied the reduction or even avoidance of software tax.
State-of-the-art (SOTA) solutions include NVMeDirect \cite{SSD:NVMeDirect:HotStorage-2016}, SPDK \cite{SSD:SPDK:CloudCom-2017}, Moneta-D \cite{SSD:Moneta-D:ASPLOS-2012}, BypassD \cite{SSD:BypassD:ASPLOS-2024}, etc.
Using the interfaces provided by NVMeDirect or SPDK, developers can fully bypass the OS kernel. Though, they must reprogram applications and manage data on their own.
Moneta-D and BypassD incorporate customized hardware for a part of management purposes. In particular, Moneta-D mainly makes a specialized storage device perform permission checks.
BypassD extends the IOMMU to map in-file offsets to disk blocks and check block-level access permissions.
Both of them also need to change the file system in the OS kernel.
\end{comment}



{\bf eBPF.}
The extended Berkeley Packet Filters (eBPF) \cite{ebpf} was initially developed to filter packets in networks. It now enables advanced capabilities like performance monitoring, security enforcement, and I/O optimization by allowing custom code to run efficiently in the OS kernel.
A notable feature of eBPF is its ability to transparently alter the execution path of a kernel function and run alternate codes that have been verified~\cite{ebpf,280870}.
This makes a seamless transition to a new execution path for a particular syscall, without requiring any change to applications.
eBPF also supports a memory buffer shared at the boundary of kernel and user spaces.
We will exploit these features to exempt the heavy tax imposed in the software stack on database logging I/Os.


\iffalse
\fi